\documentclass{article}
\usepackage{iclr2027_conference,times}
\usepackage{amsmath,amssymb}
\usepackage{booktabs}
\usepackage{enumitem}
\usepackage{graphicx}
\usepackage{hyperref}
\usepackage{url}

\iclrfinalcopy

\title{Embedding Agentic Traces for Efficient Evaluation}

\author{Anonymous\thanks{Draft skeleton; author block TBD (Helivan, Jataware, Johns Hopkins).}\\
Helivan Research
}

\newcommand{\trubric}{\textsc{trubric}}

% Notation follows arXiv:2605.07096: systems f \in F, benchmark Q* = {q_1,...,q_M},
% response-level scorer s, benchmark score y, references (f_i, y_i), budget m << M.
% The embedding function g is reserved for Section 3.

\begin{document}
\maketitle

\begin{abstract}
Efficient benchmarking estimates a system's full-benchmark performance from a small number of task executions.
Standard approaches rely solely on the binary correctness of executions from a collection of already-evaluated reference systems.
An agentic system, however, produces a rich trace for each execution that can reveal substantially more about its capabilities.
Exploiting this information requires useful trace representations.
We propose representing a trace by extracting a structured, rubric-guided summary of it with an LLM and embedding the summary with a standard text encoder.
We show that these representations preserve correctness- and task-related signal while minimizing the presence of system identity-, model family-, and harness-related signals that dominate raw trace embeddings.
Given this representation for each task execution from all reference systems, we construct low-dimensional representations of the systems, predict a new system's score from its position among the reference systems, and combine the resulting geometric prediction with a two-parameter logistic item-response model.
On SWE-bench Verified and Terminal-Bench 2.0, the proposed estimator reduces mean absolute error from 0.13 to 0.08 ($p < 10^{-4}$) and from 0.14 to 0.12 ($p < 10^{-3}$), respectively, when provided a single random executed task relative to an item-response model alone.
The same estimator correctly orders 88\% of system pairs versus 79\% for scores alone at this budget, and is never less accurate than its score-only counterpart at any execution budget we consider.
These results indicate that trace geometry and correctness outcomes carry complementary predictive signal and enable more accurate benchmark score estimation from limited agent executions.
\end{abstract}

\section{Introduction}
\label{sec:intro}

The dominant paradigm of generative model use is increasingly agentic: systems that plan, invoke tools, observe intermediate results, and iterate over long horizons rather than emit a single response \citep{yang2024sweagent, wang2025openhands}.
As agentic capabilities increase, so does the range of settings in which a given system might be deployed -- and, with it, the number of benchmarks on which the system must be evaluated before it can be trusted with consequential work \citep{kapoor2025hal}.
Agentic evaluation, however, is expensive: each task execution requires a properly sandboxed environment, the tasks themselves are complex and long-horizon, and a single execution routinely consumes hundreds of thousands of tokens \citep{jimenez2024swebench, tbench2026}.
At reported per-task costs of \$0.13--\$4 (an aggregate \$1.84 per rollout across 21{,}730 rollouts in \citealp{kapoor2025hal}), fully evaluating a single system costs roughly \$100--\$2{,}000 per benchmark, and the sensitivity of agentic systems to prompts, scaffolds, and underlying model versions means there is always another variant to evaluate \citep{kapoor2024agents}.
Notably, executing an agentic task produces more than a grade: file diffs, tool invocations, intermediate plans, and error messages are all created along the way.
These execution artifacts carry information about a system's capabilities well beyond the binary ``correct'' or ``not correct'' that standard evaluation pipelines retain.

Given the cost of a single full evaluation, efficient evaluation is necessary for increasing the evaluation coverage of a system.
Task-efficient (or query-efficient) evaluation -- benchmark compression and anchor-item selection \citep{vivek2024anchor, maiapolo2024tinybenchmarks}, item response theory and adaptive testing \citep{kipnis2025metabench, truong2025amortized} -- reduces the number of required executions by exploiting structure in the item-level scores of previously-evaluated reference systems.
These methods typically use only the binary correctness of each execution.
A complementary family of methods leverages additional item-level information -- such as cached responses or traces -- to further improve the modeling of system behavior: responses are embedded with an off-the-shelf embedding function, systems are represented as points in the resulting low-dimensional geometry, and benchmark prediction is cast as inference in that space \citep{duderstadt2023data, helm2025statistical, zhuang2025embedllm, helm2026query}.
Critically, these methods rely on the embedding function to capture what it means for two artifacts to be \emph{similar}.
For short natural-language responses, off-the-shelf semantic encoders fill this role well \citep{reimers2019sbert}.
Agentic traces are different: they are long, semi-structured logs whose surface form is dominated by harness formatting, tool-call syntax, and model-family idiosyncrasies \citep{sun2025idiosyncrasies}.
Indeed, we find that linear probes recover the identity of the system that generated a raw trace embedding with balanced accuracy up to $0.89$ -- while representations of what the agent actually did remain comparatively impoverished -- making off-the-shelf embedding functions poorly suited for extending geometry-based efficient evaluation to the agentic setting.

\noindent \textbf{Contributions.} We make two primary contributions.
\begin{enumerate}[itemsep=0pt, topsep=0pt, parsep=2pt]
\item \textbf{Methodological:} We introduce \trubric{} (\emph{trace rubric}), a representation that extracts a structured, rubric-guided summary of each trace with an LLM and embeds the summary with a standard text encoder.
Across seven embedding functions, linear probes show that \trubric{} representations preserve correctness- and task-related signal while collapsing system-identity decodability from $0.89$ to $0.20$ balanced accuracy -- selective attenuation of confounders rather than erasure of content.
\item \textbf{Empirical:} We demonstrate that \trubric{}-based estimators are query-efficient on SWE-bench Verified (107 systems) and Terminal-Bench 2.0 (65 systems) under a conservative leave-one-model-family-out protocol.
Blended with an item-response model, \trubric{} geometry reduces single-task MAE from 0.13 to 0.08 ($p < 10^{-4}$) and from 0.14 to 0.12 ($p < 10^{-3}$) on the two benchmarks respectively, correctly orders 88\% of system pairs versus 79\% for scores alone, and is never less accurate than its score-only counterpart at any budget.
Ablations show the gains grow with the reference library, compose with adaptive task selection, and are insensitive to the choice of embedding function.
\end{enumerate}

\subsection{Problem statement}
\label{sec:problem}

Given an agentic system $f$, we consider the problem of predicting its score on a benchmark $Q^{*} = \{q_{1}, \hdots, q_{M}\}$.
For our purposes, a system $f \in \mathcal{F}$ is a random mapping from a query space $\mathcal{Q}$ to a trace space $\mathcal{X}$: executing $f$ on query $q$ yields a trace $f(q)$, a long, semi-structured log of the system's actions, observations, and outputs.
Let $s: \mathcal{X} \to \{0, 1\}$ denote the harness's trace-level scoring function and $y: \mathcal{F} \times 2^{Q^{*}} \to [0,1]$ denote the benchmark scoring function, $y(f, Q) = \frac{1}{|Q|} \sum_{q \in Q} \mathbb{E}[s(f(q))]$.
We refer to $y(f, Q^{*})$ as $y$ when the context is clear.

As described above, full evaluation of $f$ on all queries may be infeasible.
Instead, we assume access to previously-evaluated (system, score) pairs $(f_{1}, y_{1}), \hdots, (f_{n}, y_{n})$, along with each reference system's trace and trace-level score for every query.
Such reference caches exist publicly at scale.\footnote{The SWE-bench Verified leaderboard archive contains traces and outcomes for over one hundred systems; we estimate it embodies \$50K--\$100K of agent compute contributed by submitters.}
Given this additional information, our goal is to estimate $y$ with $m \ll M$ queries.
That is, we seek to identify an estimate $\hat{y}$ that predicts the full benchmark score of $f$ from its traces (and, optionally, its trace-level scores) on a subset $Q \in 2^{Q^{*}}$ and the information available from the previously-evaluated systems.

Two aspects distinguish this setting from the short-response setting of \citet{helm2026query}.
First, the observation associated with each query is not a short response but a trace of $10^{4}$--$10^{6}$ tokens in one of many mutually incompatible harness formats.
Whether and how traces can be embedded usefully is the central question of this paper.
Second, the reference population is confounded: reference systems that share the target's underlying language model or harness permit trivially accurate prediction via score-copying.
We therefore evaluate under a leave-one-LLM-out protocol -- references sharing the target's underlying language model are excluded -- and verify robustness under a leave-one-harness-out protocol.
The query subset $Q$ may be fixed, drawn at random, or selected adaptively per target; we report all three regimes.

\section{Background \& related work}
\label{sec:related}

\paragraph{Agents, traces, and evaluations.}
Agentic benchmarks pair long-horizon tasks with automated grading harnesses.
Examples include SWE-bench and its verified subset for repository-level software issues \citep{jimenez2024swebench, openai2024verified}, executed by scaffolds such as SWE-agent \citep{yang2024sweagent} and OpenHands \citep{wang2025openhands}; Terminal-Bench for command-line tasks \citep{tbench2026}; and suites that aggregate many such settings \citep{kapoor2025hal}.
The traces these systems produce have recently become objects of study in their own right.
\citet{mehtiyev2026beyond} analyze leaderboard trajectories via symbolic action alphabets, \citet{shu2026traceprobe} document per-scaffold format fragmentation, \citet{cemri2025mast} taxonomize failure modes, \citet{zhuge2025agentjudge} propose trajectory-level judging, and \citet{xiao2025agentdiet} study trajectory compression.
This literature produces taxonomies and statistics rather than reusable vector representations.
Notably, its own attribution results imply that naive trace representations are dominated by confounders: coding-agent trajectories are attributable to scaffold and model at 85.7\% accuracy against an 11\% baseline \citep{oderinwale2026trajectories}, and model-level fingerprints survive paraphrasing and summarization \citep{sun2025idiosyncrasies, wu2026dna}.
Benchmark-integrity findings -- memorization inflating resolve rates \citep{liang2025illusion} and mislabeled gold patches \citep{yu2025utboost} -- further motivate representations that capture how a system works rather than relying on the pass/fail bit alone.

\paragraph{Efficient evaluations.}
Several lines of work aim to reduce the computational cost of benchmark evaluation.
\citet{vivek2024anchor} select ``anchor points'' from clusters of cross-model predictions, \citet{maiapolo2024tinybenchmarks} use item response theory to construct benchmark subsets of approximately 1\% of the original size, \citet{kipnis2025metabench} compress across benchmarks jointly, and \citet{truong2025amortized} amortize item calibration for adaptive testing.
These techniques consume only correctness patterns -- a few bits per probe -- which imposes a resolution floor at extreme budgets \citep{yauney2026reliable}.
Efficiency work specific to agents is nascent: \citet{ndzomga2026efficient} observe that subset selection degrades under scaffold shift, and \citet{ge2026psychometrics} extend psychometric models with task features while still consuming correctness rather than trace content.
Closest to our work is the Data Kernel Perspective Space (DKPS) lineage \citep{duderstadt2023data, helm2025statistical, acharyya2024consistent}, and in particular query-efficient evaluation from cached responses \citep{helm2026query}, which predicts a new model's benchmark score by comparing its responses on a small query subset to cached responses from previously-evaluated models.
We extend this framework to the agentic setting, where each observation is a trace rather than a short response and where the confounded reference population requires an invariant representation and a conservative reference protocol.

\paragraph{Embeddings of non-standard objects.}
General-purpose text embedding functions \citep{reimers2019sbert} fix one notion of similarity at training time, and no single embedding can represent all relevance relations \citep{weller2025limit}.
Instruction-conditioned variants \citep{su2023instructor} honor only coarse, trained-in task tags, and conditional similarity \citep{deshpande2023csts} formalizes similarity-under-a-criterion at the sentence scale.
A related family of methods embeds LLM-generated text in place of the original object: answers to an instruction \citep{peng2024inbedder}, answer vectors to question batteries \citep{benara2024qaemb}, and generated descriptions \citep{ravfogel2024description}.
Embeddings of interaction data exist for short dialogues \citep{liu2022dial2vec} and for reinforcement learning trajectories \citep{ge2025trajectory} -- the latter observing that trajectory embeddings encode agent ``style,'' precisely the confounder our setting must remove.
\trubric{} differs from these approaches in that it re-represents very long, heterogeneous logs at inference time without training, includes an explicit invariance mechanism (per-query consensus centering), and is evaluated with respect to confounder invariance rather than assuming it.

\section{qubric}
\label{sec:method}
%% TODO: port method section (Model 1 / Model 2 / centering; configuration;
%% pruning and full-trace inputs; cost). See notes/summary.tex.
%% Reserve g: X -> R^p for the embedding function, matching 2605.07096.

\section{Experiments}
\label{sec:experiments}
%% TODO: port from notes/summary.tex + figures/q100_tables.md:
%%   - requirements evidence (radar, heatmap, flat raw curve, stability)
%%   - efficient evaluation results (Tables 1-2, savings curve)
%%   - protocol robustness (leave-one-harness-out)
%% Style note: use claim-form \paragraph headers, e.g.
%% "\paragraph{Qubric-based methods outperform correctness-only estimators
%%  at low query budgets.}" as in 2605.07096.

\section{Discussion}
\label{sec:discussion}
%% TODO

\bibliography{paper_refs}
\bibliographystyle{iclr2027_conference}

\end{document}
